\documentclass[11pt, a4paper]{article}

% bfw-style.tex — gemeinsame Style-Definitionen für alle Handouts/Aufgabenhefte
% Voraussetzung: Kompilation mit xelatex (wegen fontspec)

\usepackage[ngerman]{babel}
\usepackage{geometry}
\geometry{a4paper, top=2.2cm, bottom=2.2cm, left=2.3cm, right=2.3cm}

\usepackage{fontspec}
% Fallback-Kette: Source Sans Pro -> Calibri -> Segoe UI -> Latin Modern Sans
% (Latin Modern ist immer mit MiKTeX vorhanden)
\IfFontExistsTF{Source Sans Pro}{%
  \setmainfont{Source Sans Pro}[Ligatures=TeX]%
}{\IfFontExistsTF{Calibri}{%
  \setmainfont{Calibri}[Ligatures=TeX]%
}{\IfFontExistsTF{Segoe UI}{%
  \setmainfont{Segoe UI}[Ligatures=TeX]%
}{%
  \setmainfont{Latin Modern Sans}[Ligatures=TeX]%
}}}
\IfFontExistsTF{Source Code Pro}{%
  \setmonofont{Source Code Pro}[Scale=0.92]%
}{\IfFontExistsTF{Consolas}{%
  \setmonofont{Consolas}[Scale=0.92]%
}{%
  \setmonofont{Latin Modern Mono}[Scale=0.92]%
}}

\usepackage{xcolor}
% Farbpalette: hell, freundlich, modern
\definecolor{bfwPrimary}{HTML}{0EA5A4}    % Petrol/Türkis - Primär
\definecolor{bfwPrimaryDark}{HTML}{0F766E} % dunkler Petrol für Headings
\definecolor{bfwAccent}{HTML}{F59E0B}     % Amber - Akzent
\definecolor{bfwSuccess}{HTML}{10B981}    % Grün - Erfolg / Lösung
\definecolor{bfwWarn}{HTML}{F97316}       % Orange - Achtung
\definecolor{bfwInk}{HTML}{0F172A}        % fast Schwarz
\definecolor{bfwInkSoft}{HTML}{475569}    % gedämpftes Grau
\definecolor{bfwBgSoft}{HTML}{F8FAFC}     % off-white
\definecolor{bfwBgTeal}{HTML}{ECFEFF}     % helles Türkis
\definecolor{bfwBgAmber}{HTML}{FEF3C7}    % helles Gelb
\definecolor{bfwBgRose}{HTML}{FFE4E6}     % helles Rosé für Eselsbrücken

\usepackage[colorlinks=true, linkcolor=bfwPrimaryDark, urlcolor=bfwPrimaryDark]{hyperref}
\usepackage{enumitem}
\setlist[itemize]{leftmargin=*, itemsep=2pt, topsep=2pt}
\setlist[enumerate]{leftmargin=*, itemsep=2pt, topsep=2pt}

\usepackage[most]{tcolorbox}
\usepackage{booktabs}
\usepackage{tikz}
\usetikzlibrary{decorations.pathmorphing, positioning, shapes.geometric, arrows.meta}

% Merkkästchen — wichtige Definition / Kernpunkt
\newtcolorbox{merksatz}[1][]{%
  enhanced, breakable,
  colback=bfwBgTeal,
  colframe=bfwPrimary,
  boxrule=0.6pt,
  arc=4pt,
  left=10pt, right=10pt, top=8pt, bottom=8pt,
  fonttitle=\bfseries\color{bfwPrimaryDark},
  title={\faLightbulb\ Merksatz},
  #1
}

% Eselsbrücke — Mnemoniks
\newtcolorbox{eselsbruecke}[1][]{%
  enhanced, breakable,
  colback=bfwBgRose,
  colframe=bfwAccent,
  boxrule=0.6pt,
  arc=4pt,
  left=10pt, right=10pt, top=8pt, bottom=8pt,
  fonttitle=\bfseries\color{bfwWarn},
  title={Eselsbrücke},
  #1
}

% Beispiel-Box
\newtcolorbox{beispiel}[1][]{%
  enhanced, breakable,
  colback=bfwBgSoft,
  colframe=bfwInkSoft,
  boxrule=0.4pt,
  arc=3pt,
  left=10pt, right=10pt, top=8pt, bottom=8pt,
  fonttitle=\bfseries\color{bfwInk},
  title={Beispiel},
  #1
}

% Lösungs-Box (für Aufgabenhefte)
\newtcolorbox{loesung}[1][]{%
  enhanced, breakable,
  colback=bfwBgAmber!40,
  colframe=bfwSuccess,
  boxrule=0.6pt,
  arc=3pt,
  left=10pt, right=10pt, top=8pt, bottom=8pt,
  fonttitle=\bfseries\color{bfwSuccess!70!black},
  title={Lösung},
  #1
}

% Glossar-Eintrag
\newcommand{\glossareintrag}[2]{%
  \par\noindent\textbf{\color{bfwPrimaryDark}#1} \enspace #2\par\smallskip
}

% Section-Stil — etwas farbiger als Standard
\usepackage{titlesec}
\titleformat{\section}
  {\Large\bfseries\color{bfwPrimaryDark}}
  {\thesection}{0.6em}{}
\titleformat{\subsection}
  {\large\bfseries\color{bfwPrimary}}
  {\thesubsection}{0.5em}{}
\titleformat{\subsubsection}
  {\normalsize\bfseries\color{bfwInk}}
  {\thesubsubsection}{0.4em}{}

% TikZ-Stil "skizzenhaft" — etwas weicher als Rough.js, aber stabil
\tikzset{
  roughbox/.style = {
    draw=bfwPrimaryDark, thick, fill=bfwBgTeal,
    rounded corners=4pt, inner sep=6pt
  },
  roughaccent/.style = {
    draw=bfwAccent, thick, fill=bfwBgAmber,
    rounded corners=4pt, inner sep=6pt
  },
  roughline/.style = {
    draw=bfwInkSoft, thick, -{Stealth[length=2.5mm]}
  }
}

% Bullet-Symbole (bewusst kein FontAwesome — vermeidet Font-Installations-Probleme)
\newcommand{\faLightbulb}{$\bullet$}
\newcommand{\faBookmark}{$\bullet$}

% Header/Footer
\usepackage{fancyhdr}
\pagestyle{fancy}
\fancyhf{}
\renewcommand{\headrulewidth}{0pt}
\fancyfoot[L]{\small\color{bfwInkSoft}\thepage}
\fancyfoot[R]{\small\color{bfwInkSoft} BFW M-064 Beschaffung im Büromanagement}


\title{\vspace*{-1.5cm}\Huge\bfseries\color{bfwPrimaryDark}Aufgabenheft Tag~15\\[0.4em]
  \large\color{bfwInkSoft}\textnormal{Vertiefung, Übung \& Reflektion --- gemischte Wiederholung im IHK-Klausurformat}}
\author{\textsf{Büroprozesse gestalten und Arbeitsvorgänge organisieren \textperiodcentered{} M-067}\\
\textsf{\small\copyright\ Can Siebert\,\textperiodcentered\,2026}}
\date{Selbstlernmaterial \textperiodcentered{} 17.07.2026}

\begin{document}
\maketitle
\thispagestyle{empty}

\begin{merksatz}[title={\bfseries So nutzen Sie dieses Aufgabenheft}]
Dieses Heft wiederholt das \emph{gesamte Modul} (Tag 1--14). Bearbeiten Sie alle Aufgaben
zunächst \emph{ohne} in die Lösungen zu schauen und in vollständigen Sätzen. Achten Sie
auf die \emph{Operatoren} (nennen, erläutern, berechnen, beurteilen) --- sie geben den
geforderten Antwort-Tiefegrad vor. Die letzte Aufgabe ist eine \emph{komplexe
Handlungssituation}, die mehrere Themenblöcke verbindet --- genau so prüft die IHK.
Erst danach öffnen Sie den Lösungsteil ab Seite~\pageref{loesungen} und vergleichen.
\end{merksatz}

\tableofcontents
\vspace{1em}
\hrule

\section{Aufgaben}

\subsection*{Aufgabe~1 --- Raum, Ausstattung \& Umgebung (12 Punkte)}

\noindent\textbf{\color{bfwAccent}YouTube-Vorbereitung:}\enspace \href{https://www.youtube.com/watch?v=lZHAXBKssA0}{Ergonomie am Bildschirmarbeitsplatz} (\url{https://www.youtube.com/watch?v=lZHAXBKssA0})

\medskip
\begin{enumerate}[label=(\alph*)]
  \item \textbf{Nennen} Sie die sechs Teilflächen, aus denen sich der Flächenbedarf eines Arbeitsplatzes ergibt. \hfill (3~P.)
  \item \textbf{Erläutern} Sie das Grundprinzip der Ergonomie und nennen Sie die Norm für den Bürodrehstuhl. \hfill (3~P.)
  \item \textbf{Nennen} Sie die Faustregel der Steh-Sitz-Dynamik. \hfill (3~P.)
  \item \textbf{Erläutern} Sie zwei Maßnahmen, mit denen sich an einem warmen Sommertag (28~°C) das Raumklima verbessern lässt. \hfill (3~P.)
\end{enumerate}

\subsection*{Aufgabe~2 --- Arbeitsschutz \& Gesundheit (12 Punkte)}

\noindent\textbf{\color{bfwAccent}YouTube-Vorbereitung:}\enspace \href{https://www.youtube.com/watch?v=uOgrdnNIJFY}{Arbeitsschutzgesetz einfach erklärt} (\url{https://www.youtube.com/watch?v=uOgrdnNIJFY})

\medskip
\begin{enumerate}[label=(\alph*)]
  \item \textbf{Nennen} Sie die drei zentralen Rechtsgrundlagen des Arbeitsschutzes mit Abkürzung. \hfill (3~P.)
  \item \textbf{Erläutern} Sie den Unterschied zwischen \emph{Belastung} und \emph{Beanspruchung}. \hfill (3~P.)
  \item \textbf{Beschreiben} Sie Zweck und gesetzliche Grundlage der \emph{Gefährdungsbeurteilung}. \hfill (3~P.)
  \item \textbf{Beurteilen} Sie: Warum ist betriebliches Gesundheitsmanagement auch wirtschaftlich sinnvoll? \hfill (3~P.)
\end{enumerate}

\subsection*{Aufgabe~3 --- Zeit- \& Terminmanagement (12 Punkte)}

\noindent\textbf{\color{bfwAccent}YouTube-Vorbereitung:}\enspace \href{https://www.youtube.com/watch?v=dWQoldjrB1s}{Eisenhower-Prinzip \& ALPEN-Methode} (\url{https://www.youtube.com/watch?v=dWQoldjrB1s})

\medskip
\begin{enumerate}[label=(\alph*)]
  \item \textbf{Erläutern} Sie die fünf Schritte der ALPEN-Methode. \hfill (3~P.)
  \item \textbf{Ordnen} Sie folgende Aufgaben nach dem Eisenhower-Prinzip ein: (1) dringende Steuerfrist heute, (2) Vorbereitung einer Besprechung nächste Woche, (3) Routine-Werbe-E-Mail. \hfill (3~P.)
  \item \textbf{Formulieren} Sie ein Ziel nach der \emph{SMART}-Regel und erklären Sie die Buchstaben. \hfill (3~P.)
  \item \textbf{Erläutern} Sie Zweck und Funktionsweise der \emph{Wiedervorlage} in der Terminüberwachung. \hfill (3~P.)
\end{enumerate}

\subsection*{Aufgabe~4 --- Besprechung, Protokoll \& Kommunikation (14 Punkte)}

\noindent\textbf{\color{bfwAccent}YouTube-Vorbereitung:}\enspace \href{https://www.youtube.com/watch?v=0hPEBUHJw94}{Protokoll schreiben nach DIN 5008} (\url{https://www.youtube.com/watch?v=0hPEBUHJw94})

\medskip
\begin{enumerate}[label=(\alph*)]
  \item \textbf{Nennen} Sie vier Aspekte, die bei der Vorbereitung einer Besprechung zu klären sind. \hfill (4~P.)
  \item \textbf{Nennen} Sie die fünf Protokollarten und die drei Funktionen eines Protokolls. \hfill (4~P.)
  \item \textbf{Beschreiben} Sie die sprachlichen Regeln des Protokolltextes. \hfill (3~P.)
  \item \textbf{Erläutern} Sie die „vier Seiten einer Nachricht'' nach Schulz von Thun. \hfill (3~P.)
\end{enumerate}

\subsection*{Aufgabe~5 --- Post, Aufbewahrung \& Ablage (12 Punkte)}

\noindent\textbf{\color{bfwAccent}YouTube-Vorbereitung:}\enspace \href{https://www.youtube.com/watch?v=XHUau_xKmaY}{Aufbewahrungsfristen nach HGB/AO} (\url{https://www.youtube.com/watch?v=XHUau_xKmaY})

\medskip
\begin{enumerate}[label=(\alph*)]
  \item \textbf{Bringen} Sie die Arbeitsschritte des Posteingangs in die richtige Reihenfolge: Verteilen, Annehmen, Öffnen, Stempeln, Vorsortieren, Kontrollieren. \hfill (3~P.)
  \item \textbf{Nennen} Sie die Arbeitsschritte des Postausgangs (Stichworte). \hfill (3~P.)
  \item \textbf{Berechnen} Sie das Fristende einer Eingangsrechnung (Buchungsbeleg) vom 12.03.2026 und \textbf{nennen} Sie die Frist samt Fristbeginn. \hfill (3~P.)
  \item \textbf{Nennen} Sie drei Ordnungssysteme der Ablage und je ein Einsatzbeispiel. \hfill (3~P.)
\end{enumerate}

\subsection*{Aufgabe~6 --- Komplexe Handlungssituation: ein Arbeitstag bei Lindner \& Partner (18 Punkte)}

\noindent\textbf{\color{bfwAccent}YouTube-Vorbereitung:}\enspace \href{https://www.youtube.com/watch?v=_m_EtsDQ-DQ}{Büroorganisation --- Grundlagen} (\url{https://www.youtube.com/watch?v=_m_EtsDQ-DQ})

\medskip
Sie arbeiten im Sekretariat der \emph{Lindner \& Partner GmbH} (Steuerberatung, 18 Beschäftigte).
An einem Arbeitstag fallen mehrere Aufgaben aus verschiedenen Bereichen an: Eine neue
Kollegin braucht einen passenden Arbeitsplatz, um 11:00~Uhr ist eine Teambesprechung,
und mit der Mittagspost kommen ein als \emph{persönlich/vertraulich} gekennzeichnetes
Einschreiben sowie eine Eingangsrechnung.

\begin{enumerate}[label=(\alph*)]
  \item \textbf{Empfehlen} und \textbf{begründen} Sie eine Büroform für eine Tätigkeit mit konzentrierter Aktenarbeit, vertraulichen Gesprächen und gelegentlicher Teamarbeit. \hfill (4~P.)
  \item \textbf{Erläutern} Sie, wie Sie Ihren Vormittag mit dem Eisenhower-Prinzip strukturieren und die Besprechung vorbereiten (mind.\ drei Aspekte). \hfill (5~P.)
  \item \textbf{Beschreiben} Sie, welche Protokollart Sie für die gefassten Beschlüsse wählen und in welcher Sprachform Sie protokollieren. \hfill (3~P.)
  \item \textbf{Erläutern} Sie die Behandlung des vertraulichen Einschreibens und \textbf{berechnen} Sie die Aufbewahrungsfrist der Eingangsrechnung. \hfill (3~P.)
  \item \textbf{Beurteilen} Sie zwei Vorteile eines digitalen DMS und nennen Sie eine Datenschutzpflicht (DSGVO). \hfill (3~P.)
\end{enumerate}

\vspace{1em}
\noindent\textbf{Punkte gesamt: 80}

\newpage

\section{Lösungen}\label{loesungen}

\subsection*{Lösung Aufgabe~1}
\begin{loesung}
\textbf{(a)} Stellfläche, Bewegungsfläche, Funktionsfläche, Benutzerfläche am Arbeitsplatz,
Verkehrswegefläche, Besucher-/Besprechungsfläche.

\textbf{(b)} Ergonomie bedeutet, \emph{die Arbeit an den Menschen anzupassen} (nicht
umgekehrt). Der Bürodrehstuhl muss die Anforderungen der \textbf{DIN EN 1335} erfüllen
(u.\,a.\ verstellbare Sitzhöhe, Rollen, Rückenlehne).

\textbf{(c)} Faustregel der Steh-Sitz-Dynamik: ca.\ \textbf{60~\% sitzen, 30~\% stehen,
10~\% gehen}.

\textbf{(d)} Maßnahmen gegen Hitze (Beispiele): Sonnenschutz/Jalousien schließen,
früh stoß-/querlüften, Wärmequellen (Geräte) reduzieren, ggf.\ Ventilator/Klimatisierung,
Getränke bereitstellen. Laut ASR~A3.5 sollen 26~°C möglichst nicht überschritten werden.
\end{loesung}

\subsection*{Lösung Aufgabe~2}
\begin{loesung}
\textbf{(a)} Arbeitsschutzgesetz (ArbSchG), Arbeitszeitgesetz (ArbZG),
Arbeitssicherheitsgesetz (ASiG).

\textbf{(b)} \emph{Belastung} sind die von außen einwirkenden Faktoren (objektiv, für alle
gleich, z.\,B.\ Lärm). \emph{Beanspruchung} ist die individuelle Auswirkung dieser
Belastung auf die einzelne Person (abhängig von Konstitution, Erfahrung, Tagesform).

\textbf{(c)} Die \textbf{Gefährdungsbeurteilung} (\S~5 ArbSchG) verpflichtet den
Arbeitgeber, Gefahren am Arbeitsplatz systematisch zu ermitteln, zu bewerten und
Schutzmaßnahmen festzulegen sowie deren Wirksamkeit zu prüfen.

\textbf{(d)} Gesundheitsmanagement senkt Fehlzeiten und Krankenstand, erhält die
Leistungsfähigkeit und Motivation, bindet Fachkräfte und vermeidet Folgekosten von
Stress/Burnout --- es ist damit zugleich Personalfürsorge \emph{und} Wirtschaftlichkeit.
\end{loesung}

\subsection*{Lösung Aufgabe~3}
\begin{loesung}
\textbf{(a)} \textbf{A}ufgaben notieren, \textbf{L}änge (Dauer) schätzen,
\textbf{P}ufferzeiten einplanen, \textbf{E}ntscheidungen/Prioritäten setzen,
\textbf{N}achkontrolle (Unerledigtes übertragen).

\textbf{(b)} (1) Steuerfrist heute = \textbf{A} (wichtig + dringend, sofort selbst
erledigen); (2) Besprechungsvorbereitung = \textbf{B} (wichtig, nicht dringend,
terminieren); (3) Werbe-E-Mail = \textbf{Papierkorb} (weder wichtig noch dringend).

\textbf{(c)} \emph{SMART} = spezifisch, messbar, attraktiv/akzeptiert, realistisch,
terminiert. Beispiel: „Ich bearbeite bis Freitag, 16~Uhr, alle 20 offenen
Eingangsrechnungen.''

\textbf{(d)} Die \textbf{Wiedervorlage} sorgt dafür, dass Vorgänge zu einem festgelegten
Termin automatisch erneut vorgelegt werden (Wiedervorlagemappe oder Kalenderfunktion),
damit Fristen und Aufgaben nicht vergessen werden.
\end{loesung}

\subsection*{Lösung Aufgabe~4}
\begin{loesung}
\textbf{(a)} Vier Aspekte (Beispiele): Datum/Uhrzeit, Raum, Teilnehmerkreis, Zeitbedarf,
Tagesordnung (TOP), Unterlagen, Protokollführung.

\textbf{(b)} Protokollarten: \emph{Wort-, Verlaufs-, Kurz-/Ergebnis-, Beschluss-,
Gedächtnisprotokoll}. Drei Funktionen: \emph{Information, Beweis-/Dokumentationsmittel,
Organisation} (Aufgabenverfolgung).

\textbf{(c)} Der Protokolltext wird \emph{sachlich}, im \emph{Präsens} und in
\emph{indirekter Rede mit Konjunktiv I/II} verfasst; Anträge und Beschlüsse werden
\emph{wörtlich} wiedergegeben.

\textbf{(d)} Jede Nachricht hat vier Seiten: \emph{Sachebene} (Inhalt),
\emph{Selbstoffenbarung} (was der Sender über sich zeigt), \emph{Beziehungsebene}
(wie er zum Empfänger steht) und \emph{Appell} (wozu er auffordern will). Missverständnisse
entstehen oft, weil Sender und Empfänger auf unterschiedlichen Ebenen „hören''.
\end{loesung}

\subsection*{Lösung Aufgabe~5}
\begin{loesung}
\textbf{(a)} Annehmen $\rightarrow$ Vorsortieren $\rightarrow$ Öffnen $\rightarrow$
Kontrollieren $\rightarrow$ Stempeln $\rightarrow$ Verteilen (anschließend Erfassen).

\textbf{(b)} Kontrollieren $\rightarrow$ Sortieren $\rightarrow$ Falzen $\rightarrow$
Kuvertieren $\rightarrow$ Wiegen $\rightarrow$ Frankieren $\rightarrow$ Aufgeben.

\textbf{(c)} Buchungsbelege sind nach aktueller Rechtslage \textbf{8 Jahre} aufzubewahren
(Neuregelung BEG~IV; AO \S~147, HGB \S~257). Die Frist beginnt mit \emph{Schluss des
Kalenderjahres} der Entstehung. Rechnung vom 12.03.2026 $\rightarrow$ Fristbeginn
31.12.2026 $\rightarrow$ \textbf{Fristende 31.12.2034}.

\textbf{(d)} Beispiele: \emph{alphabetisch} (Kunden-/Mandantennamen), \emph{numerisch}
(Rechnungs-/Belegnummern), \emph{chronologisch} (Datum, z.\,B.\ Kontoauszüge),
\emph{geografisch} (Vertriebsgebiete), \emph{alphanumerisch} (Kombination) --- drei davon
genügen.
\end{loesung}

\subsection*{Lösung Aufgabe~6}
\begin{loesung}
\textbf{(a)} Empfehlung: \textbf{Kombibüro}. Es vereint die Vorteile von Zellen- und
Großraumbüro: Die Einzelbüros ermöglichen störungsfreie, konzentrierte Aktenarbeit und
vertrauliche Mandantengespräche, die zentrale Multifunktionszone bietet Raum für
gelegentliche Teamarbeit --- passend zur beschriebenen Mischung.

\textbf{(b)} Eisenhower: dringende Steuerfristen sofort selbst erledigen (A),
Besprechungsvorbereitung terminieren (B), Unwichtiges abblocken/delegieren (C) --- ggf.\
eine „stille Stunde'' gegen Unterbrechungen. Vorbereitung der Besprechung: Tagesordnung
erstellen, Raum und Teilnehmerkreis klären, Einladung mit Frist und Unterlagen versenden,
Protokollführung festlegen.

\textbf{(c)} Für die gefassten Beschlüsse eignet sich das \textbf{Beschlussprotokoll}.
Protokolliert wird \emph{sachlich, im Präsens und in indirekter Rede (Konjunktiv I/II)};
Beschlüsse werden wörtlich festgehalten.

\textbf{(d)} Das mit \emph{persönlich/vertraulich} gekennzeichnete Einschreiben wird
\textbf{nicht geöffnet}, sondern (quittiert/erfasst) ungeöffnet an die Empfängerin
weitergeleitet. Die Eingangsrechnung als Buchungsbeleg ist \textbf{8 Jahre}
aufzubewahren (Neuregelung BEG~IV; Fristbeginn Schluss des Kalenderjahres).

\textbf{(e)} DMS-Vorteile (zwei): schneller Zugriff/Volltextsuche, gleichzeitige Nutzung,
Platz- und Kostenersparnis, revisionssichere Archivierung (GoBD). Datenschutzpflicht:
technisch-organisatorische Maßnahmen (TOM) wie Zugriffsbeschränkung/Rollenkonzept,
Verschlüsselung oder ein Löschkonzept zum Schutz personenbezogener Mandantendaten.
\end{loesung}

\end{document}
